\chapter{Of Public Order and Forbidden Speech}
\label{art:publicorder}

\articlenote{In which is set down what a citizen may not say, may not keep, and may not make, that the silo's order and the silo's peace may both be kept whole.}

\section{The Forbidden Word}
\label{sec:publicorder:forbidden-word}

No citizen may speak, write, or otherwise give voice to a wish, a request, or a suggestion that a citizen be permitted outside the silo, save the citizen's own request made directly and finally under \artref{art:cleaning}, which this Pact answers in full and does not treat as an offense.

\begin{clauses}
\item No citizen may teach, suggest, or give another citizen cause to believe that the world above is fit, in whole or in part, for the living of a life, whether by word, by writing, or by any other means.
\item A citizen who hears another citizen give voice to the wish named above, and does not report it to the Sheriff or to Judicial without delay, is answerable alongside the citizen who gave it voice, under \artref{art:crimes}.
\end{clauses}

\section{Rumor, Despair, and the Spread of Discontent}
\label{sec:publicorder:rumor-despair-and-the-spread-of-discontent}

No citizen may spread, as though it were known fact, a rumor of the silo's coming failure, of a coming shortage this Pact and its Departments have not themselves declared, or of any other matter likely to move citizens toward panic, hoarding, or unrest.

\begin{clauses}
\item No citizen may gather other citizens, beyond the ordinary gathering of a floor, a trade, or a household, or an assembly held under \artref{art:assembly}, for a purpose this Pact does not name, without leave of the Sheriff, sought and given beforehand.
\item This section does not forbid a citizen's own grief, doubt, or complaint, spoken to another citizen in the ordinary way of living; it forbids only the giving of that grief, doubt, or complaint the shape of a fact this Pact has not itself declared.
\item No citizen may give out, as a thing known, that the view shown upon the wallscreen under \artref{art:founding}, or the report of the outward sensors, is other than what the sensors gather: that it is made, that it is altered, that it is withheld, or that it is false. A citizen who doubts the view carries the doubt to the Head of Information Technology under \artref{art:it}, or to Judicial, and to no third citizen; and the Department answers the doubt in writing within one cycle, and enters that it was carried and answered.
\item Where the wallscreen or an outward sensor fails, darkens, or shows what it has not shown before, Information Technology under \artref{art:it} puts out the fact of the failure by the stations of official word; and until that is done no citizen may give out a cause for it. A citizen may say that the screen is dark, for it is. A citizen may not say why.
\end{clauses}

\section{False Teaching of the Founding and the Past}
\label{sec:publicorder:false-teaching-of-the-founding-and-the-past}

No citizen shall invent detail of the founding, the silo's building, or the time before beyond what \artref{art:founding} records, and teach that invention to another as though it were known fact.

\begin{clauses}
\item A citizen who studies or speculates upon the founding or the time before, privately and without teaching the invention to another as fact, does not offend under this section; a citizen who teaches a child, in a place of schooling under \artref{art:schooling} or otherwise, any invented detail as fact, offends gravely, whatever the citizen's own belief in the invention's truth.
\item This section binds every citizen equally, of every office, and no office excuses the teaching of invented history as fact under this section.
\end{clauses}

\section{Relics of the Time Before}
\label{sec:publicorder:relics-of-the-time-before}

A relic is any object, writing, device, or likeness made before the silo's sealing, or bearing upon its making, its history, or the world above, save such an object as this Pact or Judicial has sanctioned for a citizen's keeping.

\begin{clauses}
\item A citizen who finds or comes to hold a relic shall surrender it to the Sheriff without delay, who shall carry it to Judicial for classification; Judicial shall enter every relic so surrendered in its archive under \artref{art:judicial}, together with its finding, its classification, and its citizen. A citizen who surrenders a relic during a Forgiveness Holiday declared under \artref{art:mayor} does so under the grace that Article provides, and not under fear of the ordinary consequence of having kept it.
\item Judicial shall classify a relic by the danger it poses to the silo's order. A relic bearing upon the founding, the time before, or the world above in a way that could teach a citizen to doubt what \artref{art:founding} records is classified among the gravest, and is destroyed or sealed beyond any citizen's reach, at Judicial's own discretion, once its examination under \artref{art:it} is complete. A relic posing lesser danger may be classified otherwise, and a citizen who finds such a relic may petition Judicial to sanction its keeping.
\item A relic sanctioned for a citizen's keeping under this section remains the silo's to reclassify, and a sanction once given may be withdrawn by a later ruling of the Judge under \artref{art:judicial}.
\end{clauses}

\section{Restricted Instruments and Devices}
\label{sec:publicorder:restricted-instruments-and-devices}

No citizen may make, keep, or use any lens or glass that magnifies beyond the ordinary power of reading spectacles, any device that sends or receives word without a wire strung and licensed under \artref{art:it}, any fine instrument for the working of metal, wire, or glass, or any device made to carry a citizen's own body between floors, save under a license issued by Information Technology under \artref{art:it} or \artref{art:mechanical}.

\begin{clauses}
\item An instrument or device named in this section, found in a citizen's keeping without license, is a relic under \secrefhere{sec:publicorder:relics-of-the-time-before}, and is classified among the gravest unless Judicial finds cause, upon examination under \artref{art:it}, to classify it otherwise.
\end{clauses}

\section{The Apparatus of the Silo's Keeping}
\label{sec:publicorder:apparatus-of-the-silos-keeping}

Information Technology under \artref{art:it} and Mechanical under \artref{art:mechanical} keep throughout the silo such apparatus as the watching of the silo's vital works requires --- the sensing of the air and the water, the reading of the generator's load, the carrying of official word by wire, and such further instruments as either Department's charge requires --- and the working of that apparatus is no part of ordinary citizenship.

\begin{clauses}
\item No citizen may cover, disable, remove, open, or otherwise interfere with any such apparatus, whether or not the citizen knows its purpose; a citizen who finds such an instrument out of its place, or damaged, shall report it to Information Technology or to the Sheriff without delay, and shall not attempt its repair.
\item Interference with such apparatus is answerable under \artref{art:crimes}; where the interference endangers the works of water, air, or power under \artref{art:mechanical}, or the outward sensors under \artref{art:founding}, it is among the gravest offenses as those Articles already provide.
\item Neither Department is required to set down for any citizen where such apparatus is kept or how it is read, that being a matter of each Department's own charge and not of ordinary citizenship.
\end{clauses}

\section{The Obstruction of an Officer}
\label{sec:publicorder:obstruction-of-an-officer}

No citizen shall hinder, delay, or obstruct the Mayor, the Judge, the Sheriff, an Officer of Judicial, a Deputy or Peacekeeper, a Head of Department, or any citizen carrying an office this Pact names, in the discharge of that office, whether by deed, by the withholding of what the office may lawfully require, or by the raising of other citizens against it.

\begin{clauses}
\item A citizen so obstructing is answerable under \artref{art:crimes}; where the obstruction is of an Officer carrying out a sentence, or of Mechanical in the keeping of the works under \artref{art:mechanical}, it is answered by the mines upon a first offense.
\item A citizen who refuses an officer's requirement on the ground that it is not lawful may say so, and shall be heard on it by Judicial under \artref{art:judicial} afterward; but the refusal is at the citizen's own risk, and is not excused where Judicial finds the requirement lawful.
\item The referral of a matter to Judicial that \artref{art:mayor}, \artref{art:mechanical}, \artref{art:mines}, or \artref{art:labor} allows, made in good faith, is not obstruction, nor is the refusal \artref{art:mayor} allows to a citizen ordered to break this Pact.
\end{clauses}

\section{The Classification of a Relic}
\label{sec:publicorder:classification-of-a-relic}

A relic surrendered under \secrefhere{sec:publicorder:relics-of-the-time-before} is classified at a sitting, upon a report, and the class is entered before anything else is done with it.

\begin{clauses}
\item Judicial keeps a relic roll, apart from the roll of evidence under \artref{art:judicial}, in which every relic surrendered or found is entered with its number, the day, the citizen who surrendered it or the place it was found, the report of Information Technology under \artref{art:it}, the class given, and what was done with it; the relic roll is not open to citizens save as this section provides.
\item The Judge classifies a relic at a posted sitting under \artref{art:judicial}, with an Officer of Judicial and, where the Judge asks it, the relic examiner of Information Technology present; the citizen who surrendered the relic may attend and be heard as to how the relic came to the citizen, and no more, and the Judge enters what the citizen says.
\item A relic is classified in one of three classes. The \emph{red} class is the gravest under \secrefhere{sec:publicorder:relics-of-the-time-before}: a relic that could teach a citizen to doubt what \artref{art:founding} records, or that is a device of the kind \secrefhere{sec:publicorder:restricted-instruments-and-devices} names, save where the Judge finds cause under that Section to classify it otherwise; a red relic is destroyed or sealed as \secrefhere{sec:publicorder:relics-of-the-time-before} provides, and no petition for its sanction is heard. The \emph{second} class is a relic of no such danger that is nonetheless the silo's to keep: a tool, a fitting, a material the works or the farms can use, or a thing the Judge finds it unwise to let pass among citizens; a second-class relic is given to Mechanical for reclamation or kept in Judicial's store, and is not sanctioned. The \emph{third} class is a relic of no danger the Judge can find, and may be sanctioned for a citizen's keeping under \secrefhere{sec:publicorder:sanction-of-a-relic}.
\item The class is entered in the relic roll with the Judge's reason in a line, and the citizen who surrendered the relic is told the class, and, where the relic is red, that it will not be seen again; a citizen is not told the report of Information Technology beyond what the Judge chooses to read.
\item A relic once classified is reclassified only by the Judge at a further sitting, upon a further report or a ruling under \secref{art:judicial}{sec:judicial:rulings-where-this-pact-is-silent}, and the reclassification is entered beside the first; a relic classified red is not reclassified.
\end{clauses}

\section{The Sanction of a Relic}
\label{sec:publicorder:sanction-of-a-relic}

A relic of the third class may be a citizen's to keep, by a sanction that is the silo's to give and to take back.

\begin{clauses}
\item A citizen may petition the Judge, under \artref{art:judicial}, for the sanction of a third-class relic the citizen surrendered, or of one in Judicial's store the citizen asks for; the petition names the relic by its number and the use the citizen would put it to, and the Judge grants, limits, or refuses it at a posted sitting, entering the reason.
\item A sanction granted is entered in the relic roll and on the citizen's record under \artref{art:records}, names the citizen and the relic, and is marked upon the relic itself by the archivist where the relic will bear a mark; a sanction is for the citizen named and no other, and a sanctioned relic passes to another citizen only by a further sanction, or upon the citizen's death to the household by the Judge's leave entered.
\item A sanctioned relic is shown to the Sheriff or an Officer upon their requiring it, with the citizen's occupant card; a relic found in a citizen's keeping that is not so marked and entered is an unsanctioned relic under \secrefhere{sec:publicorder:relics-of-the-time-before}, whatever the citizen says of it.
\item A sanction is withdrawn by the Judge under \secrefhere{sec:publicorder:relics-of-the-time-before} at any sitting, for cause entered or for none, and the relic is surrendered to the Sheriff within one day of the withdrawal's being told to the citizen; a citizen who does not surrender it is answerable under \artref{art:crimes} as for the keeping of an unsanctioned relic.
\item The Judge reads at each ordinary sitting of the Assembly under \artref{art:assembly} the count of relics classified in the year, by class, and the count of sanctions given and withdrawn, without the names of citizens or the description of any relic.
\end{clauses}

\section{Writing, Likeness, and the Wall}
\label{sec:publicorder:writing-likeness-and-the-wall}

What a citizen may write and draw, and where, is set down here, that the freedom of it be known as well as its limit.

\begin{clauses}
\item A citizen may write, draw, or make a likeness of any thing within the silo, of any citizen, or of any matter this Pact does not forbid, and may keep, give, or sell it under \artref{art:commerce}; no license or sanction is required for a writing or a likeness a citizen makes, and none is a relic by reason of its making.
\item No citizen may make a likeness of the world above, save as the wallscreen itself shows it, nor write of the world above as a place of any character but that \artref{art:founding} records; a likeness or writing so made is a writing forbidden under \secrefhere{sec:publicorder:false-teaching-of-the-founding-and-the-past} and is surrendered as a relic under \secrefhere{sec:publicorder:relics-of-the-time-before}, and the citizen who made it is answerable under \secrefhere{sec:publicorder:forbidden-word} or \secrefhere{sec:publicorder:false-teaching-of-the-founding-and-the-past} as the Judge finds.
\item No citizen may write or draw upon a wall, a door, a landing, or any common surface of the silo save the notice-board of the landing under \artref{art:sheriff} and the posting-places this Pact names; a writing so made is removed by the Sheriff, and the citizen who made it is answerable under \artref{art:crimes}, and the wall is the silo's to keep clean.
\item A likeness of the silo's own structure, its floors, its stair, or its works, made beyond what the schooling of \artref{art:schooling} teaches, is shown to the Head of Mechanical before it is kept, given, or sold, and the Head may keep it for the Department's own use, entering why, or return it; a likeness of the works or the workings is not made at all save by a citizen of Mechanical or Mines for the Department's own use.
\end{clauses}
